\documentclass[11pt]{article}
\usepackage[T1]{fontenc}
\usepackage{lmodern}
\usepackage[margin=1in]{geometry}
\usepackage{microtype}
\usepackage{booktabs}
\usepackage{tabularx}
\usepackage{array}
\usepackage{amsmath}
\usepackage[hidelinks]{hyperref}
\setlength{\parskip}{0.5em}
\setlength{\parindent}{0pt}
\newcolumntype{L}{>{\raggedright\arraybackslash}X}
\newcommand{\best}[1]{\textbf{#1}}

\title{Can a Data-Rich Market Teach a Data-Poor One?\\[0.3em]
\large Transfer Forecasting for New-Market Demand}
\author{Duy Nguyen}
\date{November 8, 2026 \quad$\cdot$\quad Working draft, results v2 (exploratory, five rounds)}

\begin{document}
\maketitle

\begin{center}
\begin{minipage}{0.9\textwidth}\small\itshape
Written in the Amazon narrative style: prose first, data in tables, open questions stated plainly.
Terms are explained the first time they appear, so a reader new to forecasting can follow.
\end{minipage}
\end{center}

\section{Purpose}

When a retailer opens in a new country, it must plan inventory before that country has much sales
history. The obvious idea is to borrow the history of a market it already knows well, for example using
US demand to plan Japan. This document tests that idea on public data and asks two questions: does
borrowing a data-rich market's history improve forecasts in a data-poor market, and \textbf{when does
that transfer fail?} The aim is not to ship a model. It is to learn, with honest measurement, where
borrowing helps and where it misleads, so the next experiment is better aimed.

\textbf{The answer so far, in five rounds.} Borrowing the UK's history directly did not beat a simple
four-week average (Round 1). The error came from estimating Germany's \emph{size}, not its seasonal
\emph{timing} (Round 2). A model from macroeconomics, which separates Germany's lasting level from one-off
orders, gave the best result on the first test (Round 3). On a full year, its lead shrank to a few points and
it failed right after Christmas (Round 4). Telling it to react faster at that known turning point fixed most
of that failure (Round 5). The final model beats the simple average by 2--5 points at every history length
and is close to unbiased, but the gain is not yet statistically confirmed.

\section{Background}

\textbf{Demand} here means units sold per week. Strictly, sales are only a stand-in for demand: when a
product is out of stock, customers who wanted it can't buy it, so recorded sales undercount true demand
(\emph{censored demand}). This dataset has no stockout records, so the study uses sales and says so.

Forecasting is framed as \textbf{supervised regression}. The label is the number of units sold in a
future week. The inputs are facts known before that week, such as last week's units, the average of the
last four weeks and the week of the year. Turning a time series into a table of ``what happened 1, 2, 3
weeks ago'' is called building \textbf{lag features}. The one rule that keeps this honest is no
\textbf{leakage}: an input may never use information from on or after the week being predicted.

Amazon has published several forecasting models: DeepAR (an LSTM-based recurrent network trained as one
\emph{global model} across many products), MQ-RNN and MQ-CNN (which forecast many weeks ahead at once and
output ranges such as P50 and P90), MQ-Transformer, and Chronos (a pretrained time-series foundation
model). This study deliberately starts below all of them, because a complex model is only worth its
cost if it beats simple ones.

\section{Tenets}

\begin{enumerate}
\item \textbf{Beat a baseline before claiming anything.} A \emph{baseline} is the simplest reasonable
forecast, such as ``next week equals last week''. A model that can't beat it isn't helping, however
clever it is.
\item \textbf{Decide the test before seeing results.} The setup, metric and models were written down
before any model ran.
\item \textbf{Report error and direction together.} A model can score well by always guessing low, so
accuracy is never reported without bias.
\item \textbf{Treat small differences as noise until shown otherwise.} Every comparison carries an
uncertainty range.
\end{enumerate}

\section{The stand-in problem and the data}

Public data doesn't contain the same products sold in the US and Japan, so the study uses a stand-in with
the same shape: \textbf{UCI Online Retail II}, the transactions of a UK online gift wholesaler from
December 2009 to December 2011. After removing cancellations, returns and non-product codes, it has about
1.04 million sales lines.

\begin{center}
\begin{tabularx}{\textwidth}{L r r l}
\toprule
Market & Orders (Dec 2009--Dec 2011) & Orders per month & Role \\
\midrule
United Kingdom & 36,184 & $\sim$1,450 & Source (data-rich) \\
Germany & 753 & $\sim$30 & Target (data-poor) \\
France, Ireland, Netherlands, Spain, Belgium & 581 down to 143 each & 6--24 & Peer markets \\
Japan & 33 in total & --- & Too few to use \\
\bottomrule
\end{tabularx}
\end{center}

Germany sold 2,168 products, and 2,166 of them were also sold in the UK, so there is plenty to borrow.
Germany's weekly units are very lumpy: in 2011 they ranged from 305 to 7,015, because a single wholesale
order can triple a week's total. The UK, with 250 to 600 orders a week, is smooth.

\textbf{A known limitation:} the UK and Germany share Christmas, so this setup tests \emph{data
scarcity} more than \emph{calendar differences}. US to Japan would add a harder problem, since peaks such
as Golden Week, Obon and New Year fall in different weeks.

\section{How the test works}

The question is posed from early September 2011, before the holiday season. Each model forecasts
Germany's weekly units for the 13 weeks from September 5 to November 28, 2011, from 1 to 4 weeks ahead.
To mimic a new market, each model is shown only the most recent \textbf{1, 3, 6 or 12 months} of German
history. UK and peer history is never cut. With 1 to 6 months, Germany's own data contains no Christmas;
with 12 months, it contains one.

Accuracy is measured by \textbf{WAPE} (weighted absolute percentage error): total absolute error divided
by total actual units. Lower is better, and it stays stable when some weeks have very few sales.
\textbf{Bias} is the total forecast minus total actual, as a percentage: positive means over-forecasting.
Uncertainty comes from a \textbf{bootstrap}, which resamples the 13 test weeks many times to see how much
the score could move by chance.

Rounds 1--3 use those 13 autumn weeks. Rounds 4--5 repeat the same test for \textbf{every week of 2011}
(48 target weeks, January 3 to November 28), so each model faces quiet months, the autumn rise and the drop
after Christmas. Models are then compared with a \textbf{paired bootstrap}, which resamples the same weeks
for both models each time, so the comparison isn't blurred by which weeks happened to be drawn.

\begin{center}
\begin{tabularx}{\textwidth}{l l L}
\toprule
Model & Uses & Idea \\
\midrule
B0 Naive & Germany only & Next week = last week \\
B1 Moving average & Germany only & Next week = average of the last 4 weeks \\
B2 LightGBM & Germany only & Gradient-boosted decision trees on lag and calendar features \\
T1 Scaled UK & Germany + UK & UK's pattern, scaled by the recent Germany/UK ratio \\
T1b Top-down (author's idea) & Germany + UK & Germany's level from its whole window, timing from the UK \\
T2 Pooled LightGBM & Germany + UK & One model trained on both countries, with country as a feature \\
T1c UK lift & Germany + UK & Germany's level $\times$ the UK's seasonal rise last year \\
T3 Peer lift (author's idea) & Germany + peers & Germany's level $\times$ the peer markets' seasonal rise last year \\
T4 Signal extraction (author's model) & Germany + peers + UK & Germany's level filtered from one-off orders
(Kalman filter), $\times$ the UK's timing \\
T5 Turning point (author's fix) & Germany + peers + UK & T4, but reacting faster from mid-December
through January \\
\bottomrule
\end{tabularx}
\end{center}

\section{Results}

\subsection*{Round 1: does borrowing the UK's history help?}

WAPE (\%), lower is better, by months of German history:

\begin{center}
\begin{tabular}{l r r r r r}
\toprule
Model & 1 mo & 3 mo & 6 mo & 12 mo & Bias at 3 mo \\
\midrule
B0 Naive & 70.9 & 70.9 & 70.9 & 70.9 & $-2\%$ \\
\best{B1 4-week average} & \best{52.3} & \best{52.3} & \best{52.3} & 52.3 & $-6\%$ \\
B2 LightGBM, Germany only & can't train & 71.0 & 70.5 & 58.6 & $-2\%$ \\
T1 Scaled UK & 61.1 & 61.1 & 61.1 & 61.1 & $+20\%$ \\
T1b Top-down & 61.1 & 60.1 & 60.4 & 57.3 & $+16\%$ \\
\best{T2 Pooled LightGBM} & 72.5 & 63.5 & 60.2 & \best{52.2} & $+23\%$ \\
\bottomrule
\end{tabular}
\end{center}

The prediction going in was that borrowing would help most when Germany had little data. \textbf{It did
not.} The simple four-week average beat every borrowing model at 1, 3 and 6 months. The borrowing models
over-forecast by 15--40\%. The pooled model improved steadily as German data grew, from 72.5 to 52.2, and
tied the simple average at 12 months, with its bias falling from $+42\%$ to $+1\%$. With enough German
rows, it learned Germany's own level instead of copying the UK's. The Germany-only LightGBM couldn't
train on one month and was weak at three to six, a reminder that machine learning needs data.

\subsection*{Round 2: borrow only the size of the rise, from peers instead of the UK}

The over-forecasts raised a question: should the size of Germany's autumn rise come from smaller peer
markets rather than the large UK? T1c and T3 use the same formula, so only the source of the rise
changes.

\begin{center}
\begin{tabular}{l r r r r r}
\toprule
Model & 1 mo & 3 mo & 6 mo & 12 mo & Bias at 3 mo \\
\midrule
B1 4-week average & 52.3 & 52.3 & 52.3 & 52.3 & $-6\%$ \\
T1c UK lift & 51.3 & 53.3 & 53.3 & 52.3 & $+15\%$ \\
T3 Peer lift & \best{50.6} & \best{51.3} & \best{48.6} & \best{47.0} & $\mathbf{-38\%}$ \\
\bottomrule
\end{tabular}
\end{center}

Autumn rise, measured as September--November weekly sales divided by June--August:

\begin{center}
\begin{tabular}{l r r}
\toprule
 & 2010 & 2011 (the test) \\
\midrule
UK & 1.78$\times$ & 1.68$\times$ \\
Peers combined & 0.88$\times$ & --- \\
Germany & 1.43$\times$ & \best{1.59$\times$} \\
\bottomrule
\end{tabular}
\end{center}

T3 has the lowest WAPE at every history length, but \textbf{for the wrong reason}. The peers barely rose
in autumn 2010 (they dipped to 0.88$\times$), so T3 predicted almost no rise and under-forecast by
14--38\%. It scored well because Germany's weeks are lumpy: a forecast that stays low misses the spikes
but sits close to the many quiet weeks. Germany's real rise (1.59$\times$) was in fact close to the UK's
(1.68$\times$). That reverses the explanation offered after Round 1: the UK's rise was roughly the right
size, and the earlier over-forecasts came from how those models estimated Germany's level. The bootstrap
90\% ranges still overlapped for every model (at 3 months: B1 39--67\%, T1c 39--74\%, T3 39--61\%), so
none of these differences is settled.

\subsection*{Round 3: separate Germany's lasting level from one-off orders}

Round 2 located the problem in Germany's level, so Round 3 asks how to estimate that level from lumpy
weeks. The author's idea came from macroeconomics. In the permanent-income model of consumption, a
household's income mixes a lasting part with one-off windfalls, and it should spend based on the lasting
part. Germany's weekly sales have the same structure: a slowly moving level, plus one-off wholesale orders.
The author wrote this as a \textbf{local level model}:
\[
d_t = L_t + e_t, \qquad L_t = L_{t-1} + u_t,
\]
where $d_t$ is the week's units, $L_t$ the lasting level, $e_t$ a one-off shock with variance $\sigma_e^2$,
and $u_t$ a lasting change with variance $\sigma_u^2$. Each week, the estimate of the level moves by a
fraction $k$ of the surprise:
\[
\hat L_t = \hat L_{t-1} + k_t\,(d_t - \hat L_{t-1}).
\]
If weeks are mostly one-off noise ($\sigma_e^2 \gg \sigma_u^2$), $k$ is near 0 and a big week is ignored;
if they are mostly lasting change, $k$ is near 1. The author's first formula, $k = \sigma_u^2/(\sigma_u^2 +
\sigma_e^2)$, is a special case of the exact answer, the \textbf{Kalman filter}, which also tracks how unsure
the current estimate is ($P$):
\[
P_{t|t-1} = P_{t-1} + \sigma_u^2, \qquad
k_t = \frac{P_{t|t-1}}{P_{t|t-1} + \sigma_e^2}, \qquad
P_t = (1 - k_t)\,P_{t|t-1}.
\]

Borrowing enters in two places. A new market has too few weeks to learn its own noise ratio
$q = \sigma_u^2/\sigma_e^2$, so $q$ is fitted by maximum likelihood on the five peer markets, using only
data before each forecast date. The filter also starts from a borrowed guess, the average peer market's
weekly units, with wide uncertainty, so Germany's own data takes over as it arrives. The forecast, T4, is
the filtered level $\times$ the UK's seasonal timing from the year before.

\begin{center}
\begin{tabular}{l r r r r r}
\toprule
Model (13 autumn weeks) & 1 mo & 3 mo & 6 mo & 12 mo & Bias at 3 mo \\
\midrule
B1 4-week average & 52.3 & 52.3 & 52.3 & 52.3 & $-6\%$ \\
T1b Top-down & 61.1 & 60.1 & 60.4 & 57.3 & $+16\%$ \\
T4 level only (no seasonality) & 50.0 & 48.9 & 48.6 & 47.1 & $-27\%$ \\
\best{T4 level $\times$ UK timing} & \best{47.3} & \best{44.9} & \best{44.2} & \best{43.6} & $-10\%$ \\
\bottomrule
\end{tabular}
\end{center}

Fitting the model gave a clear answer about how much to smooth: \textbf{a lot}. The fitted $q$ was about
$10^{-7}$, meaning week-to-week movement in these small markets is almost all one-off noise. By the end of
Germany's window the weight on each new week was only 0.13 (1 month of history), 0.06 (3), 0.03 (6) and
0.02 (12). T4 was the first model that was both the most accurate and close to unbiased. It did not depend
on a lucky starting guess: with the borrowed prior halved or doubled, WAPE stayed at 44--48 for 3 to 12
months of history. Only with one month, where the prior dominates, did a doubled prior hurt (62.0). The
bootstrap range at 3 months (33--58) still overlapped B1's (39--67).

\subsection*{Round 4: the full-year test}

Thirteen weeks in one season is a small test, so Round 4 rolls the forecast date through all of 2011. The
LightGBM models are left out because they were slow and not competitive in Round 1.

\begin{center}
\begin{tabular}{l r r r r r}
\toprule
Model (48 weeks of 2011) & 1 mo & 3 mo & 6 mo & 12 mo & Bias at 3 mo \\
\midrule
B0 Naive & 65.9 & 65.9 & 65.9 & 65.9 & $-5\%$ \\
B1 4-week average & 51.4 & 51.4 & 51.4 & 51.4 & $-5\%$ \\
T1b Top-down & 56.4 & 52.2 & 52.5 & 51.9 & $+4\%$ \\
T1c UK lift & 52.0 & \best{48.5} & 50.2 & 48.8 & $+4\%$ \\
T3 Peer lift & 69.4 & 76.6 & 68.7 & 68.3 & $+12\%$ \\
T4 level only & \best{47.1} & 49.6 & 50.8 & \best{47.0} & $-7\%$ \\
T4 level $\times$ UK timing & 47.6 & 50.0 & 52.7 & 48.8 & $-1\%$ \\
\bottomrule
\end{tabular}
\end{center}

\textbf{The Round 3 lead mostly did not hold.} Over the full year, T4 beat the four-week average by 1--4
points at 1, 3 and 12 months, with almost no bias, and lost by 1 point at 6 months. In the paired
bootstrap every interval included zero. The peer-lift model (T3), which had looked best in Round 2,
collapsed to 68--77: its win had been specific to one season.

Splitting the year by quarter showed where T4 fails (table in Round 5, 3 months of history). It was best
in the holiday quarter (39.9 against B1's 51.3) and worst right after Christmas (71.0 against 42.9), when it
over-forecast January to March by 51\%. The author's diagnosis: the heavily smoothed level is slow to
notice the post-Christmas drop, just as it is slow to notice the autumn rise. The smoothing that correctly
ignores one-off orders also ignores real turns.

\subsection*{Round 5: react faster at a known turning point}

The author's fix: at times when real change is known to happen, let the filter move faster, and keep it
slow everywhere else. T5 multiplies the variance of lasting change by a factor $M$ in a turning window,
December 15 to January 31:
\[
\sigma_{u,t}^2 = \begin{cases} M\,\sigma_u^2 & \text{if week } t \text{ is in the turning window,}\\
\sigma_u^2 & \text{otherwise.}\end{cases}
\]
A larger $\sigma_u^2$ raises $P$, and so raises the weight $k$ on new weeks, only in that window. The
window comes from the calendar, which is known in advance, so it uses no test data. $M = 100$ was set as
the main value before the run (it moves $q$ from about $10^{-7}$ toward $10^{-5}$, still heavy smoothing);
$M = 10$ and $M = 1000$ are checks.

\begin{center}
\begin{tabular}{l r r r r r r}
\toprule
Model (48 weeks of 2011) & 1 mo & 3 mo & 6 mo & 12 mo & Bias 3 mo & Bias 12 mo \\
\midrule
B1 4-week average & 51.4 & 51.4 & 51.4 & 51.4 & $-5\%$ & $-5\%$ \\
T4 ($M = 1$) & 47.6 & 50.0 & 52.7 & 48.8 & $-1\%$ & $+3\%$ \\
T5, $M = 10$ & 47.6 & 49.8 & 52.2 & 48.4 & $-1\%$ & $+2\%$ \\
\best{T5, $M = 100$ (main)} & \best{47.6} & \best{48.2} & \best{49.3} & \best{46.6} & $-4\%$ & $-5\%$ \\
T5, $M = 1000$ & 48.1 & 46.1 & 45.9 & 45.6 & $-9\%$ & $-13\%$ \\
\bottomrule
\end{tabular}
\end{center}

\begin{center}
\begin{tabular}{l r r r r r}
\toprule
WAPE by quarter, 3 months of history & Q1 & Q2 & Q3 & Q4 & Bias in Q1 \\
\midrule
B1 4-week average & \best{42.9} & 57.5 & 52.0 & 51.3 & $+2\%$ \\
T4 & 71.0 & 42.1 & 53.0 & 39.9 & $+51\%$ \\
T5, $M = 100$ & 61.7 & 42.1 & 53.0 & 39.9 & $+36\%$ \\
T5, $M = 1000$ & 50.8 & 42.0 & 53.0 & 39.9 & $+10\%$ \\
\bottomrule
\end{tabular}
\end{center}

\textbf{The fix works where it should and nowhere else.} T5 ($M = 100$) cuts the first-quarter error from
71.0 to 61.7 and leaves the other quarters unchanged. It now beats the four-week average at every history
length, by 2--5 points, including 6 months, where T4 had lost. The first quarter is still its weakest:
even reacting faster, the level starts January at the December peak and needs several weeks to fall.

\textbf{Reacting harder is not simply better.} $M = 1000$ has the lowest WAPE, about 46, but under-forecasts
by 9--13\%. With 6--12 months of history, January stays inside Germany's window all year. A hard nudge
locks the level to the post-Christmas low, and the heavy smoothing afterwards is too slow to lift it back.
This is the Round 2 trap again: WAPE rewards a cautious forecast. So $M = 100$ stays the main result.

\subsection*{How sure are we?}

More than after Round 2, but not enough to claim a win. The paired bootstrap of WAPE(T5) $-$ WAPE(B1),
over the 48 weeks of 2011, gives:

\begin{center}
\begin{tabular}{l r r r}
\toprule
German history & Mean difference (points) & 90\% interval & Share of resamples T5 wins \\
\midrule
1 month & $-3.7$ & $-8.7$ to $+1.5$ & 88\% \\
3 months & $-3.2$ & $-9.8$ to $+3.1$ & 79\% \\
6 months & $-2.1$ & $-9.5$ to $+5.2$ & 68\% \\
12 months & $-5.0$ & $-11.9$ to $+2.2$ & 87\% \\
\bottomrule
\end{tabular}
\end{center}

T5 wins in most resamples at every history length, but every interval still includes zero. One year of
one lumpy series cannot confirm a gain of a few points.

\section{What we learned}

\begin{enumerate}
\item \textbf{Simple is hard to beat when the target data is thin and noisy.} After five rounds, the best
model beats a four-week average by only 2--5 points.
\item \textbf{Borrowing a pattern also borrows a level.} Most of the transfer error came from estimating
Germany's size, not its seasonal shape. The models that worked take only the \emph{timing} from the UK and
estimate the level from Germany.
\item \textbf{``Similar size'' isn't ``similar seasonality''.} Small peer markets looked like better
comparisons for Germany, but the large UK predicted Germany's rise better. The peer markets were more
useful for something else: learning how noisy a small market is.
\item \textbf{Structure from economics told us how much to trust new data.} The permanent-versus-transitory
model, fitted on the peers, said ``smooth a lot''. It also showed where smoothing fails: at real turning
points. Telling the model \emph{when} real change is expected, from the calendar, fixed most of that.
\item \textbf{A short test can crown the wrong winner.} T3 won 13 autumn weeks and collapsed over the
year; T4's 8-point lead shrank to about 2.
\item \textbf{A good score can hide a bad forecast.} Twice, WAPE favored a model that under-forecast:
T3 by up to 38\% in Round 2, and T5 with $M = 1000$ by 13\% in Round 5. At a retailer, that means
stockouts, so bias belongs next to accuracy.
\item \textbf{These echo a lesson from the author's PSB 2027 paper,} where gains seen in training did not
carry over to held-out data: a method has to be judged where it will be used, not where it looks best.
\end{enumerate}

\section{Next steps}

\begin{enumerate}
\item Expect the January drop instead of only reacting to it: add a seasonal factor for the level, so the
forecast falls on schedule after Christmas.
\item Forecast the top 20 German products, where data is thinnest and borrowing should matter most.
\item Combine sources, weighting the UK and peer rises by how well each predicted Germany in past seasons.
\item Cap or flag single very large orders, and report results with and without them.
\item Use T5's filtered level as an input to the LightGBM models, combining the economic structure with
machine learning.
\item Test neural models against T5. Amazon's Deep State Space Models (Rangapuram et al., 2018) are the
neural version of this filter: a recurrent network sets the filter's parameters for each series, which is
a natural next step from T5.
\item Add a pair with a different holiday calendar to test the calendar-shift problem directly.
\end{enumerate}

\section{FAQ}

\textbf{Why not just use a neural network?} On 106 weeks of one noisy series, simple rules and small
statistical models are strong, cheap and easy to inspect. A neural model has to earn its place by beating
them, which is next step 6.

\textbf{What does a macroeconomics model have to do with machine learning?} Both learn parameters from
data. The local level model has two (the noise variances), learned by maximum likelihood, the same
principle behind training a neural network on a loss. Its advantage here is structure: with very little
German data, a model that knows ``demand is a level plus noise'' needs far fewer examples than one that
has to discover that.

\textbf{Why weekly and not monthly?} Monthly totals are smoother and easier to forecast, but operations
such as staffing and replenishment run week by week. The top-down model (T1b) tries to get both: a
monthly level and a weekly split.

\textbf{Why is the error so high?} Germany averages about 30 orders a month, and one wholesale order can
triple a week. No model can foresee which week a large customer orders, so the goal is to be less wrong
than the alternatives, not to reach zero.

\textbf{Does this say US to Japan would fail?} No. It says that in this data, borrowing the UK's numbers
directly helped less than expected, and borrowing worked best for timing and for learning how noisy a new
market is. US to Japan adds calendar differences that this pair doesn't test; T5's turning window would
have to move to Japan's own turning points, such as New Year and Golden Week.

\textbf{Isn't the turning window a form of tuning on the test?} The window's dates come from the calendar,
which a forecaster knows in advance, and $M = 100$ was fixed before the run. $M = 1000$ scored better on
WAPE, and it was not adopted. A stricter check would choose $M$ on 2010 and test it on 2011.

\textbf{Open decision.} T5 with $M = 1000$ has about 2 points lower WAPE than $M = 100$ but is 9--13\%
low; $M = 100$ is less than 5\% low. Which should a retailer ship? This document leaves that open: it
depends on the relative cost of a stockout and of excess inventory.

\section*{Appendix: reproducibility}
\sloppy

Code and outputs are in the project folder: \texttt{src/explore.py} (data cleaning and country summary),
\texttt{src/baselines.py} (Round 1), \texttt{src/peer\_lift.py} (Round 2), \texttt{src/t4\_signal.py} and
\texttt{src/t4\_sensitivity.py} (Round 3), \texttt{src/rolling\_2011.py} (Round 4),
\texttt{src/t5\_turning.py} (Round 5), and \texttt{results/} (predictions, tables and
\texttt{RESULTS\_v1.md}). The full derivation is in \texttt{docs/DERIVATION\_T4.md}. Data: UCI Machine
Learning Repository, Online Retail II (dataset 502), downloaded November 7, 2026. The research questions,
the top-down and peer-lift ideas, the local level model and its first gain formula, the Round 4 diagnosis
and the turning-point fix, and the interpretation are the author's; the code was built with Claude Code.

\end{document}
